\subsection{Архитектурные подходы}

При создании сложных автономных миссий важно соблюдать правильную архитектуру программы. Это обеспечивает надёжность и предсказуемость поведения дрона.

\begin{infobox}[Ключевые принципы]
\begin{enumerate}
    \item \textbf{Событийно-ориентированное программирование:}
    Логика переходов строится внутри функции \texttt{callback(event)}. Дрон сам сообщает нам, когда он завершил предыдущее действие (например, \texttt{Ev.POINT\_REACHED}). Не пытайтесь угадать время полёта с помощью таймеров — используйте события.
    
    \item \textbf{Асинхронность:}
    Никогда не используйте \texttt{sleep()} или циклы \texttt{while} для ожидания. Это заблокирует автопилот. Используйте \texttt{Timer.callLater} и \texttt{Timer.new}.
    
    \item \textbf{Безопасность (Failsafe):}
    Всегда обрабатывайте аварийные события:
    \begin{itemize}
        \item \texttt{Ev.SHOCK} (Удар) — немедленно выключайте двигатели (\texttt{DISARM}).
        \item \texttt{Ev.LOW\_VOLTAGE2} (Критический разряд) — выполняйте посадку.
    \end{itemize}
\end{enumerate}
\end{infobox}

\subsubsection{Пример структуры программы}
Типовая программа состоит из:
\begin{enumerate}
    \item Объявления переменных и настроек.
    \item Функций-обработчиков для каждого этапа полёта (взлёт, полёт по точкам, посадка).
    \item Таблицы переходов (State Machine).
    \item Функции \texttt{callback}, которая переключает состояния.
\end{enumerate}
